\section*{Chapitre \ref{ch:qm:linear-time-series} --- Séries temporelles linéaires}

\begin{solution}{exo:qm:linear-time-series:1}
\omterm{def:qm:stochastic-differential-equations:ou}{Demi-vie} $\ln(1/2)/\ln0.98 = 34.3$ jours~; $\rho(20) = 0.98^{20} = 0.67$.
\end{solution}

\begin{solution}{exo:qm:linear-time-series:2}
Non~: la racine de $1 + 2z$ vaut $-1/2$, à l'intérieur du cercle unité. Le processus MA(1) de coefficient $1/2$ et de variance d'innovation quatre fois plus grande, $X_t = \eta_t +
\frac12\eta_{t-1}$, a les mêmes autocovariances ($\rho(1) = \vartheta/(1 + \vartheta^2) = 0.4$ dans les deux cas) et est inversible.
\end{solution}

\begin{solution}{exo:qm:linear-time-series:3}
$w_0 = 1$, $w_1 = -d = -0.3$, $w_2 = w_1(1 - d)/2 = -0.105$, $w_3 = w_2(2 - d)/3 = -0.0595$.
\end{solution}

\begin{solution}{exo:qm:linear-time-series:4}
$-(1 + 3 \times 0.99)/1\,000 = -0.0040$~: en moyenne, l'estimation vaut 0.986 lorsque la vraie valeur est 0.99. La \omterm{def:qm:stochastic-differential-equations:ou}{demi-vie} pour 0.99 vaut 69.0 jours~; un analyste naïf qui
estime 0.99 devrait corriger sa valeur à 0.994, soit une \omterm{def:qm:stochastic-differential-equations:ou}{demi-vie} de 114.6 jours. Un biais de 0.4~\% sur $\phi$ est un biais de 66~\% sur la \omterm{def:qm:stochastic-differential-equations:ou}{demi-vie}.
\end{solution}

\begin{solution}{exo:qm:linear-time-series:5}
Multiplier $X_t = \phi_1X_{t-1} + \phi_2X_{t-2} + \varepsilon_t$ (centré) par $X_{t-h}$, $h \ge 1$, et prendre les espérances~: $\gamma(h) = \phi_1\gamma(h - 1) + \phi_2\gamma(h - 2)$~;
diviser par $\gamma(0)$. Pour $h = 1$, $\rho(1) = \phi_1 + \phi_2\rho(1)$, donc $\rho(1) = 0.5/1.3 = 0.385$~; puis $\rho(2) = 0.5 \times 0.385 - 0.3 = -0.108$.
\end{solution}

\begin{solution}{exo:qm:linear-time-series:6}
$T = 100$~: $-2.86154 - 0.028903 - 0.000423 - 0.000040 = -2.891$. $T = 12\,574$~: $-2.862$.
\end{solution}

\begin{solution}{exo:qm:linear-time-series:7}
\texttt{spurious()}~: $|t| > 1.96$ dans 89~\% des 5\,000 régressions en niveaux, avec un $R^2$ médian de 0.17~; en régressant les différences, dans 5.0~\%, le
taux nominal. La différenciation élimine le problème de la tendance stochastique commune lorsque les séries ne sont pas cointégrées (chapitre~20).
\end{solution}

\begin{solution}{exo:qm:linear-time-series:8}
La statistique $t$ du niveau retardé dans cette régression est le $\tau$ de Dickey--Fuller, dont les valeurs critiques à 1~\% et 5~\% sont $-3.43$ et $-2.86$~: $-2.79$
n'est pas significatif à 5~\%. La \omterm{def:qm:stochastic-differential-equations:ou}{demi-vie} de 552 jours est aussi biaisée à la baisse (739 après la correction de Kendall) et très incertaine (une fourchette de 10 à 90~\% d'environ
350 à 970 jours si la vraie valeur était 739). Une opération qui attend des années le \omterm{def:qm:stochastic-differential-equations:ou}{retour à la moyenne} doit être dimensionnée pour la possibilité qu'il n'y en ait aucun.
\end{solution}

\begin{solution}{pb:qm:linear-time-series:1}
\textbf{1.} Moyenne 84.5 pb~; étendue de $-241$ à 291 pb~; 25 pb le 22 septembre 2026.
\textbf{2.} 0.9987~; le corrélogramme décroît presque linéairement, encore à 0.64 après 250 jours, signature d'une \omterm{def:qm:linear-time-series:unitroot}{racine unitaire} ou quasi unitaire.
\textbf{3.} 4.6 pb~; 0.015, 0.015, $-0.027$.
\textbf{4.} AR(10), tous les coefficients étant inférieurs à 0.035 en valeur absolue~: détectable sur 12\,574 jours, sans valeur pour la prévision.
\textbf{5.} Approximativement un \omterm{def:qm:linear-time-series:stationary}{bruit blanc}~: $\Delta X_t \approx \varepsilon_t$ avec $\psi_j$ proche de zéro pour $j \ge 1$.
\textbf{6.} $\hat\phi = 0.99875$, \omterm{def:qm:estimation:estimator}{erreur type} 0.00045, \omterm{def:qm:stochastic-differential-equations:ou}{demi-vie} de 552 jours de bourse (2.2 ans).
\textbf{7.} $t = -2.79$, et une $p$-valeur unilatérale normale de 0.003.
\textbf{8.} $(1 + 3\hat\phi)/n = 0.00032$, donc $\phi = 0.99906$ et une \omterm{def:qm:stochastic-differential-equations:ou}{demi-vie} de 739 jours (2.9 ans).
\textbf{9.} Avec 0, 5, 10 et 20 retards~: $-2.79$, $-2.74$, $-3.13$, $-3.55$, face aux valeurs critiques $-3.43$ (1~\%) et $-2.86$ (5~\%). Non rejetée à 5~\% avec
peu de retards, rejetée à 5~\% avec dix retards et à 1~\% avec vingt retards~; le critère AIC retient 34 retards et $-3.14$.
\textbf{10.} $\tau = -2.01$~: aucun élément contre une \omterm{def:qm:linear-time-series:unitroot}{racine unitaire}.
\textbf{11.} Médiane de 576 jours, fourchette de 10 à 90~\% de 352 à 969.
\textbf{12.} Dans 57~\% des historiques.
\textbf{13.} $H = 0.34$ pour les variations, $d = -0.16 \pm 0.06$~: surdifférenciée, donc le niveau a $d \approx 0.84$, soit un
processus à \omterm{def:qm:stochastic-differential-equations:ou}{retour à la moyenne} lent et hyperbolique.
\textbf{14.} Parce que la vraie mémoire n'est pas géométrique, le coefficient AR(1) estimé dépend de la longueur et de la période de l'échantillon~; et près de 1, la \omterm{def:qm:stochastic-differential-equations:ou}{demi-vie}
amplifie chaque erreur.
\textbf{15.} La distribution du chemin de \omterm{def:qm:stochastic-differential-equations:ou}{retour à la moyenne}~: jusqu'où le spread peut évoluer à l'encontre de la position et pendant combien de temps, ce qui dépend de la volatilité
et sur les queues de la mémoire, et non sur la seule \omterm{def:qm:stochastic-differential-equations:ou}{demi-vie}.
\textbf{16.} Les données ne permettent pas de trancher~: elles sont compatibles avec un \omterm{def:qm:linear-time-series:stationary}{processus stationnaire} très persistant et avec une \omterm{def:qm:linear-time-series:unitroot}{racine unitaire}.
\textbf{17.} Les 739 jours corrigés du biais, traités comme incertains entre un et quatre ans, la position étant dimensionnée pour survivre à une absence totale de retour.
\textbf{18.} Sous une \omterm{def:qm:linear-time-series:unitroot}{racine unitaire}, la statistique $t$ suit la loi de Dickey--Fuller, dont le point à 5~\% est $-2.86$, et non $-1.645$~; la vraie $p$-valeur vaut environ 0.06.
\textbf{19.} Résultat nommé~: \emph{la \omterm{def:qm:stochastic-differential-equations:ou}{demi-vie} d'un spread}~: la \omterm{def:qm:stochastic-differential-equations:ou}{demi-vie} AR(1) du spread 2s10s est de 552 jours de bourse, 739 après la correction de biais de
Kendall~; la statistique de Dickey--Fuller $-2.79$ ne rejette pas la \omterm{def:qm:linear-time-series:unitroot}{racine unitaire} à 5~\%, et les versions augmentées rejettent ou non selon le choix des retards.
\textbf{20.} Savoir si les données écartent une marche aléatoire~: pour un spread dont on espère le retour, la réponse habituelle est qu'elle n'a pas été écartée.
\end{solution}

\begin{solution}{iq:qm:linear-time-series:1}
$\ln(1/2)/\ln\phi \approx 0.69/(1 - \phi)$ périodes. Près de 1, une petite erreur sur $\hat\phi$ est une grande erreur relative sur $1 - \hat\phi$, l'estimation est biaisée à la baisse d'environ
$(1 + 3\phi)/n$, et sa distribution est asymétrique~; la \omterm{def:qm:stochastic-differential-equations:ou}{demi-vie} peut être erronée d'un facteur deux.
\iqlookfor{La formule, la sensibilité en $1/(1 - \phi)$ et le biais des petits échantillons.}
\end{solution}

\begin{solution}{iq:qm:linear-time-series:2}
Sous $\phi = 1$, le régresseur est une marche aléatoire~; l'\omterm{def:qm:estimation:estimator}{estimateur} converge au rythme $n$, et non $\sqrt n$, et sa statistique $t$ converge vers une fonctionnelle du \omterm{def:qm:brownian-motion:bm}{mouvement
brownien} (la loi de Dickey--Fuller), décalée vers la gauche~: valeur critique à 5~\% de $-2.86$ avec constante, au lieu de $-1.645$.
\iqlookfor{Asymptotique non standard et valeurs critiques.}
\end{solution}

\begin{solution}{iq:qm:linear-time-series:3}
Vérifier s'il s'agit d'une \omterm{def:qm:linear-time-series:unitroot}{régression fallacieuse}~: les deux séries sont-elles des processus à \omterm{def:qm:linear-time-series:unitroot}{racine unitaire}, et les résidus sont-ils stationnaires (un test de cointégration, chapitre~20)~? Régresser
les variations sur les variations, ou tester les résidus avec les valeurs critiques d'Engle--Granger~; examiner l'autocorrélation des résidus.
\iqlookfor{La \omterm{def:qm:linear-time-series:unitroot}{régression fallacieuse} comme première hypothèse, et la cointégration comme test.}
\end{solution}

\begin{solution}{iq:qm:linear-time-series:4}
Un AR($p$) a une \omterm{def:qm:linear-time-series:acf}{fonction d'autocorrélation partielle} qui devient nulle au-delà du retard $p$ et une \omterm{def:qm:linear-time-series:acf}{fonction d'autocorrélation} qui décroît~; un MA($q$) a l'inverse, une
\omterm{def:qm:linear-time-series:acf}{fonction d'autocorrélation} qui devient nulle au-delà du retard $q$. Un ARMA a les deux fonctions qui décroissent~; utiliser un \omterm{def:qm:linear-time-series:ic}{critère d'information} pour choisir.
\iqlookfor{Les règles de coupure et un outil de sélection de modèle.}
\end{solution}

\begin{solution}{iq:qm:linear-time-series:5}
Autocorrélations qui décroissent comme une puissance du retard et dont la somme ne converge pas~: le $d$ fractionnaire compris entre 0 et un demi. Sur les marchés~: la volatilité (valeurs absolues et carrées
rendements), volumes de transactions, signes des flux d'ordres, certains spreads et taux d'intérêt~; rarement dans les rendements eux-mêmes.
\iqlookfor{Décroissance en loi de puissance, et la volatilité comme exemple canonique.}
\end{solution}

\begin{solution}{iq:qm:linear-time-series:6}
La régression utilise la moyenne empirique, qui absorbe une partie de chaque écart, et l'\omterm{def:qm:estimation:estimator}{estimateur} est un rapport de formes quadratiques corrélées~; ces deux effets la tirent vers le bas. L'approximation de Kendall
avec moyenne estimée donne $\E[\hat\phi] - \phi \approx -(1 + 3\phi)/n$~: environ $-4/n$ près de 1, soit $-0.016$ avec une année de données quotidiennes.
\iqlookfor{Le mécanisme et l'ordre de grandeur $-(1 + 3\phi)/n$.}
\end{solution}
